\section{B1: exact Gaussian momentum--energy smoothing without a uniform minor}
\label{sec:r33-b1}

A useful anchor theorem can be proved for the actual hard-core Maxwellian preparation without conditioning on an adverse position chart. At equilibrium, positions are hard-core coupled but velocities are jointly independent Gaussian variables and independent of all positions. This distinction is preserved under exact-number conditioning, but generally not under a path-dependent tilt or exact total momentum--energy disintegration. We perform smoothing before the latter operations.

\subsection{Exact velocity factorization of the reference law}
Fix $N\ge m$ and a feasible hard-core position set $D_{\varepsilon,N}$ of positive Lebesgue measure. The canonical reference is
\[
 P^0_{\varepsilon,N}(dX,dV)=Z_X^{-1}1_{D_{\varepsilon,N}}(X)dX
               \prod_{i=1}^N(2\pi)^{-d/2}e^{-|v_i|^2/2}dv_i.
\]
Therefore a set of $m$ reserved velocities is independent of all position variables and all nonreserved velocities. No independence of anchor \emph{positions} is asserted.

Write
\[
 P_m=\sum_{i=1}^mV_i,\qquad E_m=\frac12\sum_{i=1}^m|V_i|^2,
 \qquad W_m=E_m-\frac{|P_m|^2}{2m}.
\]

\begin{theorem}[Exact regular and singular anchor law]
\label{thm:r33-b1-density}
The variables $P_m$ and $W_m$ are independent, $P_m\sim N(0,mI_d)$ and
$W_m\sim\mathrm{Gamma}(k,1)$ with $k=d(m-1)/2$. For $m>1$, the joint density is
\[
 p_m(p,e)=\frac{e^{-|p|^2/(2m)}}{(2\pi m)^{d/2}\Gamma(k)}
       \left(e-\frac{|p|^2}{2m}\right)_+^{k-1}
       e^{-(e-|p|^2/(2m))}
       1_{\{e>|p|^2/(2m)\}}.
\]
The equal-velocity set is exactly the boundary $W_m=0$ of this law. It is integrated through its gamma density, not assigned a nonexistent nonzero first-order coarea minor.
\end{theorem}

\begin{proof}
Apply an orthogonal transformation in the label index whose first row is $m^{-1/2}(1,\ldots,1)$. It sends the standard Gaussian vector in $\mathbb R^{dm}$ to independent standard Gaussian coordinates. The first $d$ coordinates are $P_m/\sqrt m$; the remaining $d(m-1)$ coordinates, denoted $Z$, are independent and satisfy $W_m=|Z|^2/2$. Polar integration gives the gamma law. The change from $(p,w)$ to $(p,e)$ has Jacobian one, proving the formula. All normalization constants follow from these two normalized factors.
\end{proof}

\subsection{A full Fourier integral, including all momentum frequencies}
Completing the Gaussian square gives
\[
 \Phi_m(u,t)=E e^{i(u\cdot P_m+tE_m)}
 =(1-it)^{-dm/2}\exp\left[-\frac{m|u|^2}{2(1-it)}\right],
\]
so
\[
 |\Phi_m(u,t)|=(1+t^2)^{-dm/4}
       \exp\left[-\frac{m|u|^2}{2(1+t^2)}\right].
\]

\begin{theorem}[An explicit anchor-size budget]
\label{thm:r33-b1-fourier}
For $s\ge0$,
\[
 \int_{\mathbb R^d\times\mathbb R}
  (1+|u|+|t|)^s|\Phi_m(u,t)|\,du\,dt<\infty
 \quad\hbox{if}\quad m>\frac{2(d+s+1)}d.
\]
Thus one fixed integer $m$ can be chosen after specifying the desired derivative budget. There is no invocation of arbitrarily many groups while keeping the number of reserved particles fixed.
\end{theorem}

\begin{proof}
Set $u=\sqrt{1+t^2}\,z$. The inner integral is bounded by a constant depending on $d,m,s$ times
$(1+t^2)^{(d+s)/2-dm/4}$. Its $t$ integral is finite precisely under the displayed sufficient strict inequality. The calculation integrates the joint conic geometry and does not falsely multiply independent stationary-phase estimates in position-coupled anchor groups.
\end{proof}

If $k>s$ for an integer $s\ge0$, the density above, extended by zero across $w=0$, belongs to $W^{s,1}(dp\,de)$. Indeed differentiation creates Gaussian polynomials in $p$ and powers down to $w^{k-1-s}$; this exponent is integrable at zero, all lower boundary traces vanish, and the tails are Gaussian/gamma integrable. This provides a second direct route to a prescribed finite smoothness order.

\subsection{Exceptional events with the correct measurability}
\begin{proposition}[Event-independent smoothing before energy conditioning]
\label{prop:r33-b1-event}
Let $\mathcal E$ depend only on all positions and nonreserved velocities. Write total additive constraints as $(P_m,E_m)+(P_{\rm rest},E_{\rm rest})$. Then under the reference law
\[
 E[1_{\mathcal E}e^{i(uP_{\rm total}+tE_{\rm total})}]
 =\Phi_m(u,t)E[1_{\mathcal E}e^{i(uP_{\rm rest}+tE_{\rm rest})}].
\]
The absolute value is at most $P(\mathcal E)|\Phi_m(u,t)|$, including when $\mathcal E$ is exceptional.
\end{proposition}

\begin{proof}
Condition on the sigma-field generated by all positions and all nonreserved velocities. The event and the rest phase are measurable there. Independence of the reserved velocities gives the exact remaining factor $\Phi_m$. The triangle inequality proves the bound.
\end{proof}

If the event depends on reserved velocities, or a path source reweights them jointly with the exterior, this proof no longer applies. A path-source theorem requires bounds on the resulting mixed amplitudes and their source derivatives. Exact momentum/energy conditioning is performed only after obtaining and dividing the appropriate densities; it is not used to justify independence that it destroys.

\subsection{Exact conditioning and source dependence}
For independent Gaussian velocities the conditional law at a regular $(p,e)$ with $e>|p|^2/(2m)$ is given by disintegrating into the fixed center-of-mass coordinate and the uniform relative sphere of radius $\sqrt{2e-|p|^2/m}$. This explicitly constructs the coarea fibre. At the boundary the relative sphere collapses and a different, degenerate fibre occurs; no uniform smooth-density assertion across all fibres is made.

A fixed compact family of linear/quadratic velocity tilts with positive precision bounded away from zero is again Gaussian. Its means, precisions, the preceding formulas and any fixed finite source derivatives are explicit. This gives a genuinely uniform finite-dimensional preparation example. It does not turn a general extensive trajectory tilt into a Gaussian velocity tilt.

\paragraph{Mechanical application still to be established.}
The path-source exact-number local theorem of the programme needs a joint amplitude estimate valid under its actual hard-sphere trajectory weights, four-source differentiated tails, and a saddle theorem on the exact regular constraint range. This chapter completely removes the false all-exterior full-rank anchor and proves a nonempty equilibrium smoothing interface, while leaving that stronger dynamic application identifiable rather than declaring it completed.
